\section{Where Is the Attack Function in Paper [3]?}

\begin{questionbox}{Question}
Paper [3] (Jedh et al.) simulates speed spoofing attacks on the ACC. Where in the paper is the mathematical representation of the simulated attacks, and how are they parameterized?
\end{questionbox}

\begin{contextbox}{Answer: Paper [3] Does Not Define a Formal Attack Function}
Paper [3] does \textbf{not} provide a closed-form mathematical function for the attack signal $a(k)$. Instead, the attacks are defined \textbf{operationally} through the CARLA simulation setup. Specifically:

\begin{enumerate}[label=\textbf{\arabic*.}]
  \item \textbf{Section 3.2} (Architecture of the Extended Adaptive Cruise Control) describes the attack conceptually: spoofed speed messages are injected into the CAN bus, and the IDS monitors for these injections.

  \item \textbf{Section 4.1} (Simulation Setup) defines the attack parameters for each scenario:
  \begin{itemize}
    \item \textbf{Scenario 2}: The ego vehicle sensor speed is spoofed to 10 km/h (a constant offset). The attack interval is fixed at a rate of 75\%, meaning spoofed messages replace real ones with probability $p = 0.75$. No IDS is used.
    \item \textbf{Scenario 3}: Same spoofing (10 km/h), same injection probability ($p = 0.75$), but now the ACC is extended with the real-time IDS. When the IDS detects the spoofing, it triggers cold braking.
  \end{itemize}

  \item \textbf{Section 4.2} (Simulation Results) shows the effects in plots but does not formalize the injection mathematically.
\end{enumerate}
\end{contextbox}

\begin{mathbox}{Reconstructing the Attack Model from Paper [3]}
From the simulation parameters, the implicit attack model is:

\begin{equation}
  \tilde{x}(k) =
  \begin{cases}
    x_{\text{spoof}} & \text{with probability } p \\
    x(k)             & \text{with probability } 1 - p
  \end{cases}
  \tag{Paper [3] Attack Model}
\end{equation}

where:
\begin{description}[leftmargin=3cm, style=nextline]
  \item[$x_{\text{spoof}}$] The constant spoofed speed value (10 km/h in Paper [3]).
  \item[$p$] Injection probability (0.75 in Paper [3]).
  \item[$x(k)$] True vehicle speed at time step $k$.
  \item[$\tilde{x}(k)$] Speed value actually received by the ACC controller.
\end{description}

This can equivalently be written using an indicator random variable $B(k) \sim \text{Bernoulli}(p)$:

\begin{equation}
  \tilde{x}(k) = (1 - B(k))\, x(k) + B(k)\, x_{\text{spoof}}
  \tag{Equivalent Form}
\end{equation}

Or as an additive attack signal:
\begin{equation}
  a(k) = B(k)\bigl(x_{\text{spoof}} - x(k)\bigr)
  \quad\text{so that}\quad
  \tilde{x}(k) = x(k) + a(k)
  \tag{Additive Form}
\end{equation}

Note that this attack signal $a(k)$ is \textbf{state-dependent} --- its magnitude depends on the current true speed $x(k)$, since it replaces the true value with a fixed spoofed value rather than adding a fixed offset.
\end{mathbox}

\begin{contextbox}{Implication for Your Paper}
Since Paper [3] defines attacks only through simulation parameters (not formal equations), your paper has the opportunity to:
\begin{enumerate}[label=(\alph*)]
  \item Provide the \textbf{formal mathematical attack model} that Paper [3] lacks.
  \item Show how this attack model integrates into the tracking control framework's error dynamics.
  \item Demonstrate via simulation that the tracking controller handles the attack, which Paper [3] addresses only through the binary cold-brake response.
\end{enumerate}

This is a clear contribution: formalizing what Paper [3] left implicit, and providing a more nuanced mitigation than emergency braking.
\end{contextbox}
